\section{Integrable Hessians and infinite mass}
\label{sec:integrable-infinite-mass}

The construction in \cref{thm:infinite-mass-witness} has a stronger linear convergence property. Its limiting current has integrable coefficients, and the smooth approximating currents converge in the coefficient $L^1$ norm. Nevertheless their determinant masses tend to infinity. Thus the obstruction is already present for an ordinary measurable matrix of integrable coefficients.

Let $X$ be a connected compact complex manifold of dimension $n$, and fix a Hermitian form $\omega$ and a Gauduchon form $\gamma$ on $X$. Let $v_+(\omega)$ denote the upper bounded-potential mass of \cref{def:bmp}, and use $\mathrm{MA}^{\mathrm{ext}}_\omega$ for the extended measure of \cref{thm:arbitrary-extended-ma}. The notation $\mathcal I(pu)$ denotes the multiplier ideal of the weight $pu$, and $\mathcal O_X$ is the structure sheaf. Put
\[
 M=\int_X\omega\wedge\gamma^{n-1}.
\]
For a real $(1,1)$-form $T$ with measurable coefficients, let $A_T$ be its Hermitian matrix relative to $\gamma$, and let $|A_T|=(A_T^2)^{1/2}$ be its positive matrix absolute value. Use the norm
\[
 \|T\|_{L^1(\gamma)}
 =\frac1n\int_X\operatorname{tr}|A_T|\,\gamma^n.
\]
For $T\geq0$ this equals $\int_XT\wedge\gamma^{n-1}$. Any smooth background norm gives the same space of integrable coefficients on compact $X$. We write $\det_\gamma T=\det A_T$, defined almost everywhere when the coefficients are integrable.

\Needspace{24\baselineskip}
\begin{proposition}\label{prop:integrable-infinite-mass}
Suppose $v_+(\omega)=+\infty$. The potential $U$ and smooth decreasing approximants $U_N$ in the proof of \cref{thm:infinite-mass-witness} satisfy the following additional conclusions. The positive current $T=\omega+\ddc U$ has integrable coefficients, and
\[
 \|\omega+\ddc U_N-T\|_{L^1(\gamma)}
 \leq2^{1-N}M.
\]
Its pointwise determinant satisfies
\[
 \mathrm{MA}^{\mathrm{ext}}_\omega(U)
 \geq (\det_\gamma T)\gamma^n,
 \qquad
 \int_X(\det_\gamma T)\gamma^n=+\infty.
\]
Moreover,
\[
 \int_X(\omega+\ddc U_N)^n\longrightarrow+\infty.
\]
The determinant measure in the statement is allowed infinite mass; equality with the non-pluripolar measure is not required.
\end{proposition}
\begin{proof}
Use the notation of that construction:
\[
 a_j=2^{-j},\quad T_j=\omega+\ddc u_j\geq0,\quad
 U_N=\sum_{j\leq N}a_ju_j\downarrow U,\quad
 \int_XT_j^n\geq2^{j(n+1)}.
\]
Every $u_j$ is smooth and nonpositive. Put
\[
 S_N=\sum_{j\leq N}a_jT_j,\qquad
 \delta_N=\sum_{j>N}a_j=2^{-N}.
\]
The smooth positive matrices $S_N$ increase pointwise. The Gauduchon identity gives
\[
 \int_XS_N\wedge\gamma^{n-1}=(1-\delta_N)M.
\]
Monotone convergence of their traces shows that the matrix series has a finite positive limit $S$ almost everywhere. Its coefficients are integrable: in a $\gamma$-unitary frame the absolute value of every coefficient is bounded by the trace. Positivity and the same trace identity give
\[
 \|S-S_N\|_{L^1(\gamma)}=\delta_N M.
\]
Thus $S_N\to S$ also in distributions. Since
\[
 \omega+\ddc U_N=S_N+\delta_N\omega
 \longrightarrow\omega+\ddc U
\]
in distributions by $L^1$ convergence of the potentials, we obtain $T=S$. The triangle inequality now gives
\[
 \|\omega+\ddc U_N-T\|_{L^1(\gamma)}
 \leq\|S_N-S\|_{L^1(\gamma)}
       +\delta_N\|\omega\|_{L^1(\gamma)}
 =2\delta_N M.
\]

We next compare the non-pluripolar measure with the pointwise determinant, using only bounded products. Fix $N$ and set $V_N=U-U_N$. The remaining partial sums show that
\[
 \delta_N\omega+\ddc V_N\geq0,\qquad V_N\leq0.
\]
For $t>0$, define
\[
 W_{N,t}=U_N+\max(V_N,-t).
\]
It is bounded, and its current is the sum of two positive factors:
\[
 \omega+\ddc W_{N,t}
 =S_N+\bigl(\delta_N\omega+\ddc\max(V_N,-t)\bigr).
\]
On $D_t=\{U>-t\}$, the inequality $U_N\leq0$ gives $V_N\geq U>-t$. Consequently $W_{N,t}=U=\max(U,-t)$ there. Locality and positivity of bounded products imply
\[
 \mathbf1_{D_t}\mathrm{MA}^{\mathrm{ext}}_\omega(U)
 \geq\mathbf1_{D_t}S_N^n.
\]
The smooth measure $S_N^n$ ignores the pluripolar pole set of $U$. Letting $t\to\infty$ proves
$\mathrm{MA}^{\mathrm{ext}}_\omega(U)\geq S_N^n$.
The positive matrices $S_N$ increase to $T$ almost everywhere, so their determinants increase to $\det_\gamma T$. Monotone convergence on every Borel set gives the asserted determinant-measure inequality. This does not presume that either limiting measure is locally finite.

For every $N\geq j$, positivity gives $S_N\geq a_jT_j$. Hence both the measurable limit and the smooth approximating current satisfy
\[
 \int_X(\det_\gamma T)\gamma^n
 \geq a_j^n\int_XT_j^n\geq2^j,
 \qquad
 \int_X(\omega+\ddc U_N)^n
 \geq a_j^n\int_XT_j^n\geq2^j.
\]
The first inequality holds for every $j$, so the integral is infinite. For the second, every $N\geq j$ obeys the bound; this proves divergence of the full sequence, without claiming that its masses are monotone.
\end{proof}

\Needspace{28\baselineskip}
\begin{corollary}\label{cor:strong-linear-infinite-mass}
If bounded mass fails, there are nonpositive potentials $w_\varepsilon\in\PSH(X,\omega)$, $0<\varepsilon<1$, such that
\[
 \omega+\ddc w_\varepsilon\geq(1-\varepsilon)\omega,
 \qquad
 \mathcal I(pw_\varepsilon)=\mathcal O_X\quad(p>0),
\]
and both the pointwise determinant integral and the extended Monge--Amp\`ere mass are infinite. As $\varepsilon\downarrow0$, the potentials converge to zero in $L^1$ and their currents converge to $\omega$ in the coefficient $L^1$ norm. Their full Gauduchon trace measures converge to $\omega\wedge\gamma^{n-1}$ in total variation.
\end{corollary}
\begin{proof}
Take $w_\varepsilon=\varepsilon U$. Its current is
\[
 T_\varepsilon=(1-\varepsilon)\omega+\varepsilon T.
\]
The lower bound and the two $L^1$ convergence assertions follow immediately. The full trace measures have integrable densities, and their total-variation distance is at most
\[
 \varepsilon\|T-\omega\|_{L^1(\gamma)}.
\]
Every multiplier ideal is trivial by the exponential integrability already proved for $U$. Since $T_\varepsilon\geq\varepsilon T$ almost everywhere, its determinant integral is at least $\varepsilon^n$ times that of $T$, hence infinite. For the extended measure, constant truncation and mixed positivity give
\[
 \mathrm{MA}^{\mathrm{ext}}_\omega(\varepsilon U)
 \geq\varepsilon^n\mathrm{MA}^{\mathrm{ext}}_\omega(U),
\]
as in \cref{cor:bmp-characterization}. Its mass is therefore infinite as well.
\end{proof}

In this construction the full current, and therefore every smooth trace of it, is absolutely continuous with respect to smooth volume. The failure of top-degree integrability does not arise from mass of that current on a pluripolar set. Linear trace control, even with total-variation convergence and trivial multiplier ideals, must still be supplemented by nonlinear estimates.
